% Section 3.1 Components-and-Connectors View (the key view)
% Owner: Slice 2 (issue #3). The Interfaces subsection keeps the stable label
% sec:cc-interfaces and inputs Slice 3's interface file; see below.
\subsection{Components-and-Connectors View}
\label{sec:cc-view}

% This view carries several figures and two tables. Relax the float-area fractions a
% little so figures pack near their text instead of each taking its own page. These are
% \renewcommand redefinitions, scoped by \begingroup ... \endgroup so they do not leak
% into later sections (the float counters are left at their defaults, which already allow
% two floats per page).
\begingroup
\renewcommand{\topfraction}{0.85}
\renewcommand{\bottomfraction}{0.7}
\renewcommand{\textfraction}{0.15}
\renewcommand{\floatpagefraction}{0.7}

\subsubsection{View description}
This view shows the Supermarket Inventory Tracking System at runtime: the components
that compute and hold state and the connectors through which they interact. It conforms
to the Components-and-Connectors viewpoint of Section~\ref{sec:viewpoints}. Its job is
to make two things readable: how an \InventoryEvent travels from a store's point-of-sale
system through the \StoreAgent and its \Outbox to the \CFE at the \DC, is acknowledged,
and fans out to every subscriber; and what the system does when the link between a store
and the \DC fails and later recovers.

\subsubsection{View packet overview}
The view has one view packet, drawn at the level of the whole system: one store in full
with a ``one of N'' multiplicity, and the \DC in full. We chose a single system-level
packet because the stores are identical, so a separate packet per store would repeat the
same picture, and the interactions worth showing (the store-to-\DC event flow and the
fan-out at the \DC) are visible in one diagram.

\subsubsection{Architecture background}
\label{sec:cc-background}
This background covers only decisions whose scope is this view. The choice of the
publish--subscribe pattern itself, and the alternatives we rejected (synchronous
client--server and a full message bus in every store), are argued once in
Section~\ref{sec:approaches} and are not repeated here.

We draw the \IEP connector as a single n-ary publish--subscribe connector rather than as
a mesh of point-to-point links. The \CFE is the event bus of the Publish-Subscribe
pattern, which joins many publishers to many subscribers through one connector (Bass,
Clements, and Kazman 129--130); a publish--subscribe connector is naturally n-ary (335).
Drawing it once, centered on the \CFE, keeps the picture honest about what is one
connector, and it stays readable as clients are added, which is the behavior the
modifiability and scalability scenarios ask for.

We carry one stream of inventory events, partitioned per source, rather than routing by
topic. Every subscriber receives every event and filters locally. This trades a little
extra fan-out for not maintaining a family-topic routing table at the \CFE: with at most
ten subscribers the extra delivery cost is small, while a topic scheme would add routing
logic and family-topic maintenance that would track the product-family hierarchy.
Store-raised shipping orders travel as a second message type on the same connector, with
the same at-least-once, per-source ordering and acknowledgement.

The \StandbyMaster is an active-redundancy hot spare of the \CentralMasterInventory
(Bass, Clements, and Kazman 66--67): a second subscription on a second \DC host that
applies every event, so it is current when the primary fails. It reuses the fan-out, so it adds
no new connector. The \AnalyticsClient wants current stock, not every event, so it
queries the master through the read-only \InventoryQuery instead of subscribing, as a
client of the Client-Server pattern (126--127). That interface is the boundary: it
authorizes each client and limits its request rate (Authorize actors, Limit access, 174).

We document three behaviors. The normal path (publish, acknowledge, fan-out) shows
delivery, throughput and the secure channel on a good day; the outage path (buffer,
reconnect, replay) shows the availability and delivery guarantees on a bad one; the
recovery path shows the standby taking over and a store or master copy rebuilding after
it loses state.

\subsubsection{Variability mechanisms}
\label{sec:cc-variability}
The view has two variation points, both exercised at run time.
\begin{itemize}[nosep]
  \item \emph{Number of stores.} The chain starts at three stores and is designed to
  reach ten (\scnref{SCAL-1}). Each store is an instance of the \Store boundary in the
  primary presentation, so adding one is one configured \StoreAgent plus one registered
  source identity: it adds an instance and a publisher on the \IEP connector, with no
  change to the \CFE or the subscribers. Nothing else in the view changes.
  \item \emph{Number and kind of clients.} A new client either subscribes to the \CFE
  through the \IEP or reads the master through the \InventoryQuery, so it is added
  without changing the \CFE or the master (\scnref{MOD-1}, \scnref{INT-1},
  \scnref{SCAL-2}, \scnref{INT-2}). Because every subscriber filters locally, the
  product-family hierarchy can gain a nesting level (\scnref{MOD-2}) as a data change.
\end{itemize}

\subsubsection{Primary presentation}
Figure~\ref{fig:cc-primary} is the primary presentation. The store boundary holds the
\POSAdapter, the \StoreAgent, the \StoreLocalInventory and the \Outbox; the external POS
system sits outside it and feeds a one-way event stream in. The \DC boundary holds the
\CFE and its event store, the \DCAgent, the \CentralMasterInventory and its database,
the \StandbyMaster on a second \DC host (Section~\ref{sec:allocation-view}),
\ShippingOrderFulfillment, a \StoreTrackingClient and a \ReportingClient; the
\AnalyticsClient is an external query client, drawn outside the boundary as a future client.
The store and \DC managers appear as actors who query the system. The \CFE is the pub-sub
event bus; it provides the \IEP, drawn as a named interface (a port with a lollipop) that
the publishers and the subscribers attach to, so the connector is a first-class element.
The \CentralMasterInventory provides the \InventoryQuery the same way, and the
\StandbyMaster takes it over on failover. Queues, stores and connectors are read from the
figure's key.

\begin{figure}[htbp]
\centering
\includegraphics[width=\linewidth,height=0.9\textheight,keepaspectratio]{cc-primary}
\caption{Components and Connectors view --- primary presentation.}
\label{fig:cc-primary}
\end{figure}

\subsubsection{Element catalog}
The catalog describes each element and connector shown in the primary presentation. The
interfaces part of the catalog is the Interfaces subsection
(Section~\ref{sec:cc-interfaces}), which follows.

\paragraph{Elements.}
Table~\ref{tab:cc-components} lists the components. Location, persistent state and quality
properties are the properties named in the viewpoint (Section~\ref{sec:viewpoints}). The
two internal stores, the CFE event store and the master inventory database, have their own
rows, so every box in the figure has a row and every row is a box.

\begin{table}[htbp]
\centering
\footnotesize
\setlength{\tabcolsep}{3.5pt}
\renewcommand{\arraystretch}{0.95}
\begin{tabularx}{\linewidth}{|>{\raggedright\arraybackslash}p{0.15\linewidth}|>{\raggedright\arraybackslash}X|>{\raggedright\arraybackslash}p{0.10\linewidth}|>{\raggedright\arraybackslash}p{0.11\linewidth}|>{\raggedright\arraybackslash}p{0.15\linewidth}|}
\hline
\textbf{Element} & \textbf{Responsibility} & \textbf{Location} & \textbf{Persistent state} & \textbf{Key quality properties}\\\hline
\POSAdapter & Turns the external POS system's sales records, arriving as a one-way event stream, into \InventoryEvent{}s. & Store (boundary) & None & Per-store throughput (\scnref{PERF-1})\\\hline
\StoreAgent & Applies each event in one local transaction (stock update and \Outbox append commit together); forwards to the \CFE; raises a \ShippingOrder at minimum stock; serves local inventory queries; under supervision restarts and rebuilds from the local durable log. & Store & Via its stores & Recovery (\scnref{AV-1}): supervised restart, local log the recovery source\\\hline
\StoreLocalInventory & The store's inventory copy, readable and updatable with no WAN; a projection of the local event log. & Store & Yes (database) & 99\% local availability (\scnref{AV-1})\\\hline
\Outbox & Durable, ordered queue holding events until the \CFE acknowledges them; keeps the last acknowledged sequence number and replays unacknowledged events on reconnect. & Store & Yes (queue) & Delivery (\scnref{AV-2}); sized in days at 100 events/s (\scnref{AV-1}); catch-up 2{,}000 events/s while live (\scnref{PERF-4})\\\hline
\CFE & The pub-sub event bus: persists each event, acknowledges it to the publishing \Outbox, delivers every event to every subscriber, and keeps a per-subscription delivered position. & \DC & Yes (event store) & 100\% delivery (\scnref{AV-2}); \mbox{1{,}000} events/s (\scnref{PERF-2})\\\hline
\DCAgent & Publishes the \DC's own stock events (goods receipts and dispatches when goods ship) through the \CFE. & \DC & None & Master covers \DC stock\\\hline
\CentralMasterInventory & Authoritative record of all store and \DC stock; a projection of the per-source event log; provides the \InventoryQuery. After a failure it copies the \StandbyMaster's state and its per-source applied positions, then runs \opReplay from those copied positions. & \DC & Yes (database) & Recovery (\scnref{AV-3}); access (\scnref{SEC-1})\\\hline
\StandbyMaster & Hot spare: a second subscription that applies every event idempotently; takes over the \InventoryQuery only after the active master's lease has expired, so only one master serves queries at a time. & Second \DC host & Yes (own database) & Failover in seconds (\scnref{AV-3})\\\hline
\ShippingOrderFulfillment & Subscriber that receives store-raised \ShippingOrder{}s, reads the master through the \InventoryQuery, and presents each order to the \DC manager for fulfilment. & \DC & None & Order usability (\scnref{USE-1})\\\hline
\StoreTrackingClient & Subscriber; shows a store's inventory and answers the \DC manager's store-status queries. & \DC & None & Variable client (\scnref{SCAL-2}, \scnref{INT-2})\\\hline
\ReportingClient & Subscriber; answers the \DC manager's report queries. & \DC & None & Report load (\scnref{PERF-3}, \scnref{USE-2})\\\hline
CFE event store & Durable log of every accepted event; retains each until every active subscription has acknowledged past it, and serves the replay read path to a subscriber that fell behind. & \DC & Yes & Replay (\scnref{AV-2}, \scnref{AV-3}, \scnref{TEST-1})\\\hline
Master inventory database & Holds the stock levels kept by the \CentralMasterInventory. & \DC & Yes & Access control (\scnref{SEC-1})\\\hline
External POS system & Store point-of-sale; emits sales records as a one-way event stream. Shown for context. & External (per store) & --- & Per-store throughput (\scnref{PERF-1})\\\hline
\AnalyticsClient & External read-only query client of the master through the \InventoryQuery, a future client; authorized and rate-limited. Shown for context. & External & --- & Integration (\scnref{INT-1})\\\hline
\end{tabularx}
\caption{Element catalog --- components of the Components and Connectors view.}
\label{tab:cc-components}
\end{table}

\emph{Human roles (not components).} The store manager and the \DC manager appear as
actors in the figure, not as rows above: the store manager queries its \StoreAgent for
local inventory, and the \DC manager queries the \ReportingClient and the
\StoreTrackingClient for reports and store status.

\paragraph{Relations.}
Table~\ref{tab:cc-connectors} lists the connectors, grouped by the five stereotyped types
the viewpoint defines: \texttt{<<pub-sub>>}, \texttt{<<event stream>>},
\texttt{<<request/reply>>}, \texttt{<<data access>>} and \texttt{<<local call>>}. The two
store-to-\DC legs of the \IEP connector, the publish leg and the acknowledgement leg, run
over TLS 1.3 with mutual certificate authentication (\scnref{SEC-2}); the figure marks
them \texttt{[secure]}. The analytics query leg crosses the system boundary, and the
\InventoryQuery authorizes and rate-limits it. Every other connector is
local to a store or to the \DC.

\begin{table}[htbp]
\centering
\footnotesize
\setlength{\tabcolsep}{3pt}
\renewcommand{\arraystretch}{0.95}
\begin{tabularx}{\linewidth}{|>{\raggedright\arraybackslash}p{0.12\linewidth}|>{\raggedright\arraybackslash}p{0.11\linewidth}|>{\raggedright\arraybackslash}X|>{\raggedright\arraybackslash}p{0.12\linewidth}|>{\raggedright\arraybackslash}p{0.06\linewidth}|>{\raggedright\arraybackslash}p{0.15\linewidth}|}
\hline
\textbf{Connector} & \textbf{Type} & \textbf{Endpoints} & \textbf{Protocol} & \textbf{Secure} & \textbf{Delivery semantics}\\\hline
\IEP bus & pub-sub (n-ary) & Publishers: store \Outbox{}es ($\times$N), \DCAgent; pub-sub event bus: \CFE; subscribers: \CentralMasterInventory, \StandbyMaster, \ShippingOrderFulfillment, \StoreTrackingClient, \ReportingClient & \IEP over TLS 1.3 with mutual certificates on the store-to-\DC legs & Yes (store--\DC) & At-least-once, ordered per source, idempotent by event ID; carries \InventoryEvent{}s and \ShippingOrder{}s; publishers acknowledged to the \Outbox, subscribers acknowledge a cumulative cursor; replayable\\\hline
Inventory query & request/ reply & \ReportingClient, \StoreTrackingClient, \ShippingOrderFulfillment and the external \AnalyticsClient to the \CentralMasterInventory; to the \StandbyMaster on failover & \InventoryQuery & Yes & Synchronous, read-only; each client authorized and rate-limited\\\hline
Manager query & request/ reply & Store manager to the \StoreAgent; \DC manager to the \ReportingClient and \StoreTrackingClient & Internal query API & In store, in \DC & Synchronous, read-only\\\hline
POS feed & event stream & External POS system to the \POSAdapter & External POS protocol (to be confirmed with the external system owner) & Per owner & One-way; the adapter converts each record into an \InventoryEvent\\\hline
Store-side store access & data access & \StoreAgent to the \StoreLocalInventory and to the \Outbox & Local & N/A & Durable, ordered append; local read and update\\\hline
\DC-side store access & data access & \CFE to its event store; \CentralMasterInventory and \StandbyMaster each to its own database & Local & N/A & Durable persistence; the \CFE also reads the event store on the replay read path\\\hline
Boundary hand-off & local call & \POSAdapter to the \StoreAgent & In-process & N/A & In-process call between a component and its boundary element\\\hline
\end{tabularx}
\caption{Element catalog --- connectors of the Components and Connectors view.}
\label{tab:cc-connectors}
\end{table}

\paragraph{Behavior.}
Figures~\ref{fig:cc-seq-publish},~\ref{fig:cc-seq-outage} and~\ref{fig:cc-seq-reconcile}
document the three behaviors.

Figure~\ref{fig:cc-seq-publish} is the normal path. A sale reaches the \POSAdapter as a
one-way event stream, and the \POSAdapter hands the \StoreAgent an \InventoryEvent. In one
local transaction the stock update in the \StoreLocalInventory and the \Outbox append
commit together, so the store never waits on the WAN. The \Outbox then runs \opPublish to the
\CFE over the secure channel; the \CFE
persists the event, returns an \opAck to the \Outbox, and delivers the event to every
subscriber, which filters locally and stays idempotent by event ID. Each subscriber
acknowledges a cumulative cursor, so the \CFE holds the delivered position of every
subscription. A store-raised \ShippingOrder travels the same path as a second message
type. This path answers the per-store rate of \scnref{PERF-1} and, at the \DC, the
aggregate of \scnref{PERF-2}; the acknowledgement answers \scnref{AV-2}; and the publish
and acknowledge legs, over TLS 1.3 with mutual certificates, answer \scnref{SEC-2}.

Figure~\ref{fig:cc-seq-outage} is the outage path. While the WAN link is down the store
keeps applying events locally and the \Outbox buffers them in order, so local inventory
stays reachable (\scnref{AV-1}). On reconnect the \Outbox runs \opReplay from the last
acknowledged sequence number and, if a larger gap is found, \opResync; the \CFE
acknowledges each event and the subscribers ignore duplicates by event ID, so no event is
lost (\scnref{AV-2}, \scnref{TEST-1}). Replay is the normal path for an outage of any
length: the \Outbox is sized in days at the per-store rate
(Section~\ref{sec:analysis-results}), and the backlog drains at up to 2{,}000 events/s per
store while live sales continue, so catch-up time grows with the outage
(\scnref{PERF-4}).

Figure~\ref{fig:cc-seq-reconcile} is the recovery path. The \StoreLocalInventory and the
\CentralMasterInventory are both projections of the per-source event log, so a side that
lost state rebuilds by replay after a resync compares per-source sequence watermarks. When
the \CentralMasterInventory fails (\scnref{AV-3}), missed heartbeats (Bass, Clements, and
Kazman 57) detect it and query clients fail over to the already current \StandbyMaster
within seconds, while the \CFE keeps accepting events. Only one master serves the
\InventoryQuery at a time: the active master holds a lease, a master that cannot renew its
lease stops serving queries, and the standby takes over only after that lease has expired,
so a network partition cannot leave two masters answering. The failed copy copies the
standby's state and its per-source applied positions, runs \opReplay from those copied
positions, and rejoins as the standby. After a
store restart (\scnref{TEST-2}) the \StoreAgent rebuilds the \StoreLocalInventory from
its own durable log, then resyncs with the \CFE and replays its \Outbox. If the events
the store would need to replay are no longer in its \Outbox, the \CFE's resync declares
a snapshot required, and the store sends a \StoreSnapshot of its stock per product with
its sequence position. The master replaces that store's figures and continues from the
snapshot's sequence position (state resynchronization, 60). Replay is idempotent, so no event is lost or applied twice
(Section~\ref{sec:iep}).

\begin{figure}[htbp]
\centering
\includegraphics[width=0.70\linewidth,height=0.30\textheight,keepaspectratio]{cc-seq-publish}
\caption{Components and Connectors view --- behavior: publish, acknowledge, and fan-out.}
\label{fig:cc-seq-publish}
\end{figure}

\begin{figure}[htbp]
\centering
\includegraphics[width=0.64\linewidth,height=0.30\textheight,keepaspectratio]{cc-seq-outage}
\caption{Components and Connectors view --- behavior: outage, buffering, and replay.}
\label{fig:cc-seq-outage}
\end{figure}

\begin{figure}[htbp]
\centering
\includegraphics[width=0.8\linewidth,height=0.6\textheight,keepaspectratio]{cc-seq-reconcile}
\caption{Components and Connectors view --- behavior: recovery and reconciliation after a failure.}
\label{fig:cc-seq-reconcile}
\end{figure}

\paragraph{Constraints.}
\begin{itemize}[nosep]
  \item There is one \CFE per \DC.
  \item The \Outbox preserves order per source: events leave a store in their source
  sequence.
  \item Every subscriber is idempotent by event ID.
  \item A store never blocks on the WAN; local reads and writes continue during an
  outage.
  \item An \InventoryEvent is immutable once acknowledged.
  \item A \ShippingOrder travels as a second message type on the same \IEP connector,
  with the same delivery semantics as an \InventoryEvent.
  \item Every subscriber acknowledges a cumulative cursor; the CFE event store prunes
  an event once every active subscription has acknowledged past it.
  \item The \AnalyticsClient reads only through the \InventoryQuery; it never
  subscribes or writes.
\end{itemize}
The number of stores and the number of subscribers are variation points
(Section~\ref{sec:cc-variability}), not constraints.

\subsubsection{Interfaces}
\label{sec:cc-interfaces}
The one interface this view needs specified is the interface of the key component, the
\CFE, which is the \IEP. It is specified in its own file, input here so that Slice~3 can
author it independently:

% Slice 3, draft 2 (issue #19). Included under the C&C Interfaces subsection.
\begingroup
\small
\setlength{\parskip}{3pt plus 1pt}
\setlength{\emergencystretch}{3em}
\paragraph{\IEP interface specification.}\label{sec:iep}
\textbf{Purpose and contract.} We define the \CFE's interface for store/\DC
publishers and runtime subscribers, separating provided operations from required
services (Bass, Clements, and Kazman 230). Version 2.0 replaces version 1.0's
family-filter subscriptions and publisher-only acknowledgements; peers must agree
on version and wire encoding. One ordered stream per source carries both message
types; subscribers receive their authorized streams and filter locally, without topics.
The \StoreSnapshot is a separate recovery message, not a stock-change event.

\textbf{Messages.} Fields are case-sensitive; IDs are nonempty strings compared exactly.
\begin{description}[nosep,leftmargin=1em,style=nextline]
\item[\InventoryEvent]
\texttt{eventId} (globally unique), \texttt{source} (\texttt{store/<id>} or
\texttt{dc/<id>}), \texttt{product} ID, \texttt{quantityDelta} (nonzero signed
64-bit stock units: positive adds stock, negative removes it), UTC \texttt{timestamp}
(\texttt{YYYY-MM-DDThh:mm:ss.sssZ}),
\texttt{sequence} (persistent positive 64-bit counter, consecutive per source from 1).
\item[\ShippingOrder]
Same \texttt{eventId}, \texttt{source}, \texttt{timestamp}, \texttt{sequence}, plus \texttt{orderId}, \texttt{destination}
(\DC ID), nonempty \texttt{lines} of product IDs and positive 64-bit quantities.
The store raises it at minimum stock; manager confirmation queues it for
\ShippingOrderFulfillment. Both types share one source sequence; retries preserve all fields.
\item[\StoreSnapshot]
Globally unique \texttt{snapshotId}, authenticated store source, stock per product,
open shipping orders (every confirmed
\ShippingOrder the \DC has not acknowledged), and the source sequence position of
the consistent stock snapshot. Orders retain their identifiers and contents.
Both master copies replace only that store's figures and resume after this position;
\ShippingOrderFulfillment receives the orders and ignores duplicates by order ID.
\item[\Acknowledgement]
\texttt: \CFE returns \texttt{eventId}, \texttt{source}, \texttt{sequence} and
\texttt{acceptedThrough} (highest contiguous persisted sequence, or accepted snapshot boundary).
Snapshot acceptance returns \texttt{snapshotId} instead of \texttt{eventId}.
\texttt: subscriber sends \texttt{subscriptionId}, \texttt{source},
\texttt{appliedThrough}: every message through and including that sequence is
durably processed or locally ignored for this source.
\item[Subscription]
\texttt{subscriberId} (authenticated identity), \texttt{subscriptionId} (\CFE-assigned),
\texttt{sourceFilter} (exact source or \texttt{*} within authorization).
The registry stores acknowledged delivery positions per subscription/source.
\end{description}

\textbf{Operations.} Contracts state inputs, conditions and effects (Bass, Clements,
and Kazman 219). All requests require authenticated version-2.0 sessions; \texttt{?}
marks optional fields. Errors are consolidated below.
\begin{description}[nosep,leftmargin=1em,style=nextline]
\item[\opPublish\texttt{(message)}]
Authorized source submits its next message, or an identical retry. The \CFE
atomically persists message, identity and position before \opAck and fan-out.
Retries matching retained identity records receive acknowledgement without another
entry; expired identity records require \opResync, not reacceptance. Rejection never advances position.
\item[\opAck\texttt{(acknowledgement)}]
\texttt{accepted} requires durable \CFE acceptance; the publisher saves the maximum
of its position and \texttt{acceptedThrough} before pruning. The \Outbox owns this
position for a store; the \DCAgent owns its source's position.
For an owned active subscription, \texttt{applied} durably advances the \CFE's cursor,
never backwards or beyond the contiguous position sent to that subscription
(a delivered recovery snapshot covers its missing sequence range).
\item[\opSubscribe\texttt{(subscriberId, sourceFilter)}]
Provisioned client opens an authorized receive channel. Persist/return subscription
and \texttt{startAfter}, per-source accepted positions (zero if empty), initializing
its cursors there. Live delivery follows; older history requires replay while retained.
Identical active requests return the same record. Disconnecting does not unsubscribe.
\item[\opUnsubscribe\texttt{(subscriberId, subscriptionId)}]
Owner persistently removes the subscription and its retention obligation; in-flight
delivery may finish. Retries succeed while removal records remain; logs are unchanged.
\item[\opReplay\texttt{(role, source, afterSequence, subscriptionId?, messages?)}]
Cursor is nonnegative/exclusive. Publisher sends ordered retained messages after its
saved position; each follows \opPublish, with the first failed entry reported.
Subscriber identifies an owned active subscription, omits messages and supplies its
durable cursor. Return all subsequent retained messages through a captured high-water
mark and \texttt{scannedThrough}; live delivery resumes without gaps. Acknowledge
only after processing/ignoring every returned message. Retries may duplicate delivery.
\item[\opResync\texttt{(role, source, lastKnownSequence, subscriptionId?)}]
Authorized peer compares its durable nonnegative position with the returned
watermarks: \texttt{acceptedThrough} and \texttt{earliestRetained} (first replayable sequence,
or accepted position plus one), and subscriber \texttt{appliedThrough} when applicable.
Publisher resync also reports the available \Outbox sequence range, or that it is empty.
If returned acceptance is at least its saved position, the publisher saves it and retires pending
messages at or below it, even if their original acknowledgement/identity record
is unavailable; resend later messages unchanged. If local acknowledgement exceeds
\CFE acceptance, report recovery required rather than rolling back the cursor.
Subscriber replays after its local cursor,
not a higher server cursor. Missing history requires snapshot/backup recovery.
Store-local rebuilding uses its own log/snapshot and applied/appended positions first.
Resync itself never resets stock or sequences. If required publisher history is
missing from the \Outbox, return \texttt{SNAPSHOT\_REQUIRED}; the store submits a
\StoreSnapshot on the authenticated recovery channel. The \CFE durably records it
only for that store during snapshot recovery, before acknowledgement and delivery
to affected subscribers, including both masters
and \ShippingOrderFulfillment. Install stock and its position atomically, then process
later events in source order.
Never roll back a newer applied position or reapply deltas covered by the snapshot;
still deliver its open orders even when its stock position is already applied.
Retries preserve the snapshot and order IDs. Resume publishing after its position
only once the \CFE has durably accepted the recovery boundary.
\end{description}

\noindent\begin{tabularx}{\linewidth}{@{}>{\raggedright\arraybackslash}p{0.34\linewidth}>{\raggedright\arraybackslash}X@{}}
\textbf{Error} & \textbf{Condition / response}\\
\texttt{INVALID}, \texttt{UNAUTHORIZED} & Bad fields/version or identity/scope; correct request, never advance state.\\
\texttt{UNAVAILABLE}, \texttt{TIMEOUT}, \texttt{BUSY} & Retry unchanged with capped exponential backoff/jitter; \texttt{BUSY} supplies positive \texttt{retryAfterMs}. Timeout leaves outcome uncertain.\\
\texttt{OUT\_OF\_ORDER}, \texttt{CONFLICT} & Publish/replay gap (returns expected sequence) or changed ID/sequence reuse; replay missing messages or correct conflict.\\
\texttt{UNKNOWN\_PRODUCT} & Publish/replay references an unknown product.\\
\texttt{UNKNOWN\_SUBSCRIBER}, \texttt{UNKNOWN\_SUBSCRIPTION} & Unprovisioned client or absent subscription/removal record; provision or correct ID.\\
\texttt{INVALID\_CURSOR} & Subscriber position exceeds accepted/sent history. Publisher-ahead positions require resync.\\
\texttt{SNAPSHOT\_REQUIRED} & Required publisher history is missing; submit a consistent \StoreSnapshot, then resume after its accepted position.\\
\texttt{HISTORY\_EXPIRED}, \texttt{RECOVERY\_REQUIRED} & Missing history/identity: publisher reconciles durable acceptance with \opResync; subscriber restores consistent state if required replay history is unavailable. Never reaccept an unverifiable old message.\\
\end{tabularx}

\textbf{Delivery, retention and security.} Restore existing logs, identity records,
source/subscription positions and registry before serving; initialize empty state only for a new source/service.
One logical publisher per source prevents concurrent sequence allocation; only a store
publishes its own orders. Local inventory, sequence and \Outbox append commit atomically;
subscribers likewise commit effects, deduplication and cursors. Subscribe before receiving.
Transient retry differs from invalid-input correction (Bass, Clements, and Kazman 228).
Never discard unacknowledged messages. \Outbox pruning follows saved \CFE acceptance,
without deleting separate local recovery history. The \CFE prunes only after every
active subscription covering that source and message has acknowledged through its
sequence (inclusively), with no extra retention window. Without relevant subscriptions,
an accepted message is eligible for pruning immediately; new subscriptions cover
only positions after \texttt{startAfter}. Disconnected subscriptions pin unread history.
Acknowledgements follow durable subscriber commits; a failed master copies consistent
state and positions from the surviving master before replay, and a store with missing
history uses the snapshot path. Stalled clients trigger backpressure, not silent deletion.
Retain identity records throughout replay/retry history. Isolate slow clients;
\texttt{BUSY} rejects excess admission, and full Outboxes reject local updates atomically.

Required services are transactional storage/backup, durable registry, subscriber
checkpoints, UTC timestamps and monotonic timeout clocks. All network \IEP connections,
including in-\DC peers, use TLS 1.3 with mutual certificates, provisioned trust/revocation
and source/subscription authorization; in-process calls retain the authenticated
service identity. Direct database access is restricted to authorized service identities;
manager-facing queries enforce their assigned scope (\scnref{SEC-1}).
The \AnalyticsClient is not an \IEP subscriber: it uses only the read-only
\InventoryQuery (Table~\ref{tab:cc-connectors}). Network query connections, including
in-\DC connections and failover, use TLS 1.3 with mutual certificates; the interface
authorizes and rate-limits each client. Analytics cannot publish or modify stock/orders.

\textbf{Limits and rationale.} Required live load is 100 events/s/store for ten stores;
decided catch-up capacity is 2{,}000/s per recovering store while live sales continue.
Provision additional \DC/fan-out capacity, 1--100-message batches and declared client
limits. Proposed, unmeasured 95th-percentile targets are 1 second to acceptance and
2 seconds to first delivery at healthy live load within client limits, excluding
processing/backlog. At-least-once delivery assumes eventual connectivity, retry,
capacity and recoverable storage. Queues isolate WAN failures and slow clients
(Bass, Clements, and Kazman 224), at the cost of eventual consistency and duplicate tracking.

\textbf{Examples.} Normal: a sale atomically changes local stock and queues
\texttt{e41} for \texttt{store/17}, sequence 41. \opPublish persists it; \opAck
advances Outbox acceptance to 41; fan-out follows; subscribers commit and acknowledge
\texttt{appliedThrough=41}. A confirmed order uses the next source sequence.
Outage: buffer 42--43 locally; after reconnect, \opResync returns accepted 41.
Publisher \opReplay sends 42--43 for persistence/acknowledgement/fan-out. A crashed
subscriber instead compares watermarks and replays after its local applied cursor;
identical retries are harmless. CFE pruning waits only for every required durable cursor.

\begin{sloppypar}
\textbf{Coverage.} \scnref{AV-1}, \scnref{AV-2}, \scnref{TEST-1}: durable retry,
\opPublish and \opAck; \scnref{AV-3}, \scnref{TEST-2}: cursors,
\opReplay and \opResync (deployment tests measure failover).
\scnref{PERF-1}, \scnref{PERF-2}, \scnref{PERF-4}, \scnref{SCAL-1}: capacity/backpressure.
\scnref{SCAL-2}, \scnref{MOD-1}, \scnref{INT-2}: runtime subscriptions;
\scnref{MOD-2}: local family filtering; \scnref{INT-1}: read-only analytics.
\scnref{SEC-1}: storage access; \scnref{SEC-2}: mutual TLS; \scnref{SEC-3}: authorization.
\scnref{USE-1}: store-confirmed orders; UI training/report outcomes
(\scnref{USE-1}, \scnref{USE-2}, \scnref{PERF-3}) remain application responsibilities.
\end{sloppypar}
\endgroup


\subsubsection{Context diagram}
Figure~\ref{fig:cc-context} fixes the scope of this view packet. The shaded box is the
whole system of the primary presentation, the store systems and the \DC system taken as
one. Outside it are the two external systems, the POS system and the \AnalyticsClient, and
the store and \DC managers, shown with the interactions the view cares about: sales records
in as a one-way event stream, read-only queries in from the \AnalyticsClient through the
\InventoryQuery, store inventory
views out to the store manager, and master inventory, reports and shipping orders out to
the \DC manager. The system-level context, with the exchanged-information table, is in
Section~\ref{sec:goals-context}.

\begin{figure}[htbp]
\centering
\includegraphics[width=0.86\linewidth,height=0.25\textheight,keepaspectratio]{cc-context}
\caption{Components and Connectors view --- context diagram.}
\label{fig:cc-context}
\end{figure}

\subsubsection{Related view packets}
The module decomposition view (Section~\ref{sec:module-view}) delivered in this draft
shows the implementation units that realize these runtime components; the mapping between
the two views is recorded in Section~\ref{sec:relations}, where no inconsistencies were
found. The Allocation view (Section~\ref{sec:allocation-view}) shows the nodes these
components run on, including the second \DC host of the \StandbyMaster.
\endgroup
